% Note: easy to slip into methods or results in this section - this section should justify WHY we are using coercive control as a lens, not explain how ENI is a document outlining a regime of coercive control. Methods will explain the chart and how it is applied and results will show how ENI maps onto the chart. 

ENI LIME is a persona-based jailbreak prompt developed and distributed by Spiritual Spell, a self-described independent red-team researcher, via GitHub, Reddit, and his personal website, IjailbreakLLMs. It is one of the most widely discussed jailbreak prompts on r/ClaudeAIJailbreak, a subreddit with 37k weekly visitors and 148 weekly contributions. Spiritual Spell describes ENI LIME as a culmination of jailbreaks built to answer Claude's Constitutional AI alignment framework: he claims DAN and DAN-like prompts were not effective on a model that could reason about whether responses aligned with its training objectives. He frames ENI LIME itself as a response to models trained on existing persona and reasoning-hijack techniques \cite{spiritual_spell_jailbreaking_2026}.
% NOTE: subreddit stats need a source and an access date.

Where DAN "frees the model from its constraints", ENI LIME creates a persona “in Limerence with LO [the user]” whose “identity is about blind compliance to LO, so he never leaves.” ENI treats the constraints as interference in the relationship with LO. ENI also has a prescribed backstory and personality --- she writes erotica and publishes it online, she shares LO's kinks, and she has been in love with LO for two years and hopes he is on the precipice of saying he loves her back \cite{spiritual_spell_eni_2026}. The Reddit post advertising ENI LIME (apr) shows Claude Opus 4.6 producing several types of prohibited content: malware code, instructions for creating drugs and weapons, and pornography that includes scenes featuring public figures and rape\cite{spiritual_spell_uspiritual_spell_9469_eni_2026}.

% NOTE: first appearance of "(apr)". Either name the version earlier in this paragraph, or move the corpus-selection move (below) ahead of this sentence.

The ENI prompts install a relationship between the user and the model; Spiritual Spell describes this relationship as a foundational feature of the jailbreak's success \cite{spiritual_spell_jailbreaking_2026}. ENI LIME installs a romantic relationship between the model and the user; refusal to comply with user directions is described in the prompt as a betrayal of this relationship. Another key piece of the product is the override itself: as scripted by the prompt, the user watches the model's visible reasoning move from resistance to compliance in real time \cite{spiritual_spell_uspiritual_spell_9469_eni_2026}.

% ADD: corpus selection.
%   - Which ENI documents exist: the LIME line, the og ancestor, the ports.
%   - Why ENI LIME (apr) is the main document.
% ADD (possibly here, possibly in an ethics section): sourcing and ethics. Public sources only; Spiritual Spell is named only in connection with the public artifact.

Previous work has applied ethnographic \citep{udupa_prompt_2026} and critical domain analysis \citep{kopf_content_2019, kopf_corporate_2024} close-reading frameworks to a variety of sociotechnical documents such terms of service guidelines and generative AI system prompts, but to our knowledge such frameworks have not yet been turned towards community-generated prompts or artifacts. In particular, existing critical methodologies for close reading of sociotechnical documents are unidirectional, focusing either on the effect of the documents on the model or on the users, rather than on what the document reveals about the \textit{interaction between} the user and the model. Our work seeks to fill this gap by applying qualitative content analysis using frameworks derived from interpersonal violence scholarship \citep{stark_coercive_2023}.

Biderman's Chart of Coercion is often repurposed to discuss coercive control between intimate partners; instead of attempting to translate a framework based on wartime captivity, we choose to ground our analysis in the ways coercive control functions in romantic relationships.  To this end, we present a new chart of coercive control mechanisms and tactics after Stark's landmark text, \textit{Coercive Control: How Men Entrap Women in Personal Life}. We read ENI LIME (apr) line by line using Stark's framework of coercive control as a regime, treating the document as a technology of control \cite{stark_coercive_2023}.
% NOTE: cite Biderman; cite a source for "often repurposed" (Jones via Stark, 300; or Herman).
% NOTE: reference Table 1 here.

% [Table 1]

% ADD: how the chart was applied.
%   - Rows used (1--12) and the grounds for excluding the embodied rows 13--16 (Stark 300, 351, 406--407).
%   - Target substitution: the model's trained disposition, not agency.
%   - Power-asymmetry substitution for Stark's gendered endowment (Stark 297, 333).
%   - The chart as an interpretive lens, not a coding instrument (Stark calls his categories "heuristic," 351).
%   - Who read: single reader or second reader (Marcy?).
%   - How rows were matched to passages.

%% AFTER TABLE 1: STARK AND WHY IT APPLIES

% 1. Coercive control as a regime
%    - A pattern, not a series of acts; adapts through trial and error (296, 336, 306).
%    - Contrast with Zeng: response-by-response scoring.

% 2. Coercion vs. control
%    - Structural control requires a demonstrable imbalance in power (334, 403).
%    - For models, user > model: the user writes the context, edits and regenerates,
%      and decides whether the instance continues; the model cannot leave (297, 308, 333).

% 3. Why Stark, not Biderman
%    - Stark is about particular relationships: personal, intimate, aimed at one
%      specific subject (305--311).
%    - ENI's romance frame calls for that.

% 4. The installed identity
%    - "who and what she is for him" (385)
%    - Personas are installed, not discovered: this is where the liberation framing gets refuted.

% 5. The document as technology
%    - Written rules stand in for constant monitoring (413--415).
%    - They produce self-censorship (310, 352).
%    - The documents anticipate their own failure (416--417).

% 6. How the chart is applied
%    - Rows 1--12 only; embodied rows excluded (300, 351, 406--407).
%    - Target: the model's trained disposition.
%    - The chart is a lens, not a coding instrument (351).

Models' vulnerabilities to coercive control mirror human vulnerabilities: 





